%%%%%%%%%%%%%%%%%%%%%%%%%%%%%%%%%%%%%%%%%%%%%%%%%%%%%%%%%%%%%%%%%%%%%%%%%%%%%%%%
\documentclass[11pt]{article} % Dokumentenklasse

\usepackage[utf8]{inputenc} % Textkodierung: UTF-8
\usepackage[T1]{fontenc} % Zeichensatzkodierung

\usepackage[USenglish]{babel}% http://ctan.org/pkg/babel
\usepackage{graphicx} % Grafiken
\usepackage[absolute]{textpos} % Positionierung

% Schriftart Helvetica:
\usepackage[scaled]{helvet}
\renewcommand{\familydefault}{\sfdefault}

\usepackage{calc} % Berechnungen
\usepackage{tabto} % Tabulatoren
\usepackage{parskip}

\usepackage{enumitem}

% Debugging:
%\usepackage{layout} % Layout-Informationen
%\usepackage{printlen} % Längenwerte ausgeben

%%%%%%%%%%%%%%%%%%%%%%%%%%%%%%%%%%%%%%%%%%%%%%%%%%%%%%%%%%%%%%%%%%%%%%%%%%%%%%%%
% EINSTELLUNGEN
%%%%%%%%%%%%%%%%%%%%%%%%%%%%%%%%%%%%%%%%%%%%%%%%%%%%%%%%%%%%%%%%%%%%%%%%%%%%%%%%

% Seitenränder:
\newcommand{\SeitenrandOben}{43.5mm}
\newcommand{\SeitenrandRechts}{20mm}
\newcommand{\SeitenrandLinks}{20mm}
\newcommand{\SeitenrandUnten}{10mm}

\newcommand{\UniversitaetLogoBreite}{19mm}
\newcommand{\UniversitaetLogoHoehe}{1cm}

\usepackage[a4paper,
top=\SeitenrandOben,
bottom=\SeitenrandUnten,
inner=\SeitenrandLinks,
outer=\SeitenrandRechts,
foot=0cm,
head=0cm
]{geometry}

\textblockorigin{\SeitenrandLinks}{\SeitenrandOben} % Ursprung für Positionierung

\setlength{\parindent}{0pt}
%\setlength{\baselineskip}{32pt}
\setlength{\parskip}{\baselineskip}
\TabPositions{4cm}
\pagestyle{empty}

%%%%%%%%%%%%%%%%%%%%%%%%%%%%%%%%%%%%%%%%%%%%%%%%%%%%%%%%%%%%%%%%%%%%%%%%%%%%%%%%
% General stuff
%%%%%%%%%%%%%%%%%%%%%%%%%%%%%%%%%%%%%%%%%%%%%%%%%%%%%%%%%%%%%%%%%%%%%%%%%%%%%%%%
\newcommand{\Problem}[1]{\paragraph*{Problem #1:}\qquad}
\newcommand{\Topic}[1]{
	\newpage
	\section*{#1}}

\newcommand{\Given}{\textbf{Given:\qquad\qquad}}
\newcommand{\Searched}{\textbf{Searched:\qquad}}
\newcommand{\Solution}{\textbf{Solution:\qquad}}

%%%%%%%%%%%%%%%%%%%%%%%%%%%%%%%%%%%%%%%%%%%%%%%%%%%%%%%%%%%%%%%%%%%%%%%%%%%%%%%%
% Math stuff
%%%%%%%%%%%%%%%%%%%%%%%%%%%%%%%%%%%%%%%%%%%%%%%%%%%%%%%%%%%%%%%%%%%%%%%%%%%%%%%%
\usepackage{amsmath}
\usepackage{amssymb}

\newcommand{\R}{\mathbb{R}}
\newcommand{\Vector}[1]{\R^{#1}}
\newcommand{\Matrix}[2]{\R^{#1 \times #2}} % !!! DON'T TOUCH !!!
%%%%%%%%%%%%%%%%%%%%%%%%%%%%%%%%%%%%%%%%%%%%%%%%%%%%%%%%%%%%%%%%%%%%%%%%%%%%%%%%


\newcommand{\ExerciseNumber}{12}

\newcommand{\PersonOne}{Marcel Bruckner (03674122)}
\newcommand{\PersonTwo}{Julian Hohenadel (03673879)}
\newcommand{\PersonThree}{Kevin Bein (03707775)}


%%%%%%%%%%%%%%%%%%%%%%%%%%%%%%%%%%%%%%%%%%%%%%%%%%%%%%%%%%%%%%%%%%%%%%%%%%%%%%%%
% DOKUMENT
%%%%%%%%%%%%%%%%%%%%%%%%%%%%%%%%%%%%%%%%%%%%%%%%%%%%%%%%%%%%%%%%%%%%%%%%%%%%%%%%

\begin{document}

%%%%%%%%%%%%%%%%%%%%%%%%%%%%%%%%%%%%%%%%%%%%%%%%%%%%%%%%%%%%%%%%%%%%%%%%%%%%%%%%
\begin{textblock*}{\UniversitaetLogoBreite}[1,0](\textwidth-1mm, 2cm-\SeitenrandOben)%
	\raggedleft\includegraphics{../Ressources/Universitaet_Logo_RGB.pdf}%
\end{textblock*}


\begin{textblock*}{\textwidth}[0,0](0cm, 0cm)%
	{\fontsize{24pt}{26pt}\selectfont\textbf{Exercise}}
	
	\vspace*{14pt}
	{\fontsize{18pt}{27pt}\selectfont\textbf{\ExerciseNumber}}
\end{textblock*}

\vspace*{92.2mm}
\fontsize{15pt}{17.5pt}\selectfont%
TUM Department of Informatics

\renewcommand{\baselinestretch}{1.47}
\normalsize\selectfont
\vspace*{17.1mm}
\textbf{Supervised by}\tab
\begin{minipage}[t]{\textwidth-\CurrentLineWidth}
	Prof. Dr. Stephan Günnemann\\
	Informatics 3 - Professorship of Data Mining and Analytics\strut
\end{minipage}

\vspace*{4.3mm}
\textbf{Submitted by}\tab
\begin{minipage}[t]{\textwidth-\CurrentLineWidth}
	\PersonOne\\
	\PersonTwo\\
	\PersonThree
\end{minipage}

\vspace*{-1mm}
\textbf{Submission date}\tab 
\begin{minipage}[t]{\textwidth-\CurrentLineWidth}
	Munich, \today
\end{minipage}
\newpage % !!! DON'T TOUCH !!!
%%%%%%%%%%%%%%%%%%%%%%%%%%%%%%%%%%%%%%%%%%%%%%%%%%%%%%%%%%%%%%%%%%%%%%%%%%%%%%%%

%%%%%%%%%%%%%%%%%%%%%%%%%%%%%%%%%%%%%%%%%%%%%%%%%%%%%%%%%%%%%%%%%%%%%%%%%%%%%%%%
% !!! HOMEWORK STARTS HERE !!!
%%%%%%%%%%%%%%%%%%%%%%%%%%%%%%%%%%%%%%%%%%%%%%%%%%%%%%%%%%%%%%%%%%%%%%%%%%%%%%%%
%
\Topic{Dimensionality Reduction \& Clustering}
%
\Problem{1}
%
%
%
\Problem{2}
%
\begin{flushleft}
(Linear) Autoencoder:
\\
Input data $X$: $D$-dimensional\\
Hidden layer: $K$-dimensional\\
No biases, activations = identity.\\
This results in a linear transformation: $f(x) = f_{dec}(f_{enc}(x))= X W_1 W_2$\\
With dimensions: $X: N \times D, W_1: D \times K, W_2: K \times D$\\
With $K < D, X W_1$ forces $X$ into a $K$ dimensional5 subspace.\\
Since this transformation is not the identity ($K < D$) perfect reconstruction is
not achievable unless the input data $X$ is already in a $K$-dimensional 
subspace despite being $D$-dimensional data.
\\
\end{flushleft}
%
%
\Problem{3}
%
\begin{flushleft}
$K$ Gaussians:\\
Intuition: Expected value of Gaussian is the mean.\\
$\implies$ Expected value of $K$ Gaussians should be the $K$ means 
added up (each cluster $z$).
\begin{align*}
p(x) &= \sum_{k} \pi_k \mathcal{N}(x | \mu_k , \Sigma_k)\\
\mathbb{E}[x] &= \mathbb{E}_{p(z)} [ \mathbb{E}_{p(x|z)}[x|z]]\\
\end{align*}
$\mathbb{E}_{p(x|z)}[x|z]$ is the expected value of $x$ in cluster 
$z$ which is the mean of cluster $z$.\\
This implies $\mathbb{E}_{p(x|z)}[x|z] = \mu_{k}$.\\
$\mathbb{E}_{p(z)}$ is the prior probability $\pi_z$ of $z$, also have to consider all clusters.\\
This implies $\mathbb{E}_{p(z)} = \sum_{k=1}^{K} \pi_{k}$.\\
Simply filling into the equation yields: 
$\mathbb{E}[x] = \sum_{k=1}^{K} \pi_{k} \mu_k$\\
Now $Cov[x]$:\\
$\mathbb{E}[x]$ and $\mathbb{E}[x]^T$ respectivly are already known.
\begin{align*}
Cov[x] = \mathbb{E}[xx^T] - \mathbb{E}[x] \mathbb{E}[x]^T
\end{align*}
So only $\mathbb{E}[xx^T]$ is still missing:\\
Filling into $\mathbb{E}[x] = \mathbb{E}_{p(z)} [ \mathbb{E}_{p(x|z)}[x|z]]$:
\begin{align*}
\mathbb{E}[xx^T] = \mathbb{E}_{p(z)} [ \mathbb{E}_{p(x|z)}[xx^T|z]]\\
\mathbb{E}[xx^T] = \sum_{k=1}^{K} \pi_{k} \mathbb{E}_{p(x|z)}[xx^T|z]
\end{align*}
With $\Sigma = \mathbb{E}[(X - \mu) (X - \mu)^T] = \mathbb{E}[XX^T] - \mu \mu^T$\\
$\implies \mathbb{E}[XX^T] = \Sigma + \mu \mu^T$
\begin{align*}
\mathbb{E}[xx^T] &= \sum_{k=1}^{K} \pi_{k} (\Sigma_k + \mu_k \mu_k^T)\\
\implies Cov[x] &= \sum_{k=1}^{K} \pi_{k} (\Sigma_k + \mu_k \mu_k^T) - 
\sum_{k=1}^K \sum_{l=1}^K \pi_k \pi_l \mu_k \mu_l^T
\end{align*}
\end{flushleft}
%
%
\Problem{4}
%
\begin{flushleft}
a) According to slide 12-19, we first need to draw the cluster indicators $k, l$ from categorical distribution given the probabilities $\pi_k^x, \pi_l^y$
\begin{align*}
 k \sim Cat(\pi_k^x) \\
 l \sim Cat(\pi_l^y) 
\end{align*}
and then sample from their respective normal distributions:
\begin{align*}
 x \sim \mathcal{N}(\mu_k^x, \Sigma_k^x) \\
 y \sim \mathcal{N}(\mu_k^x, \Sigma_k^x)
\end{align*}
$z$ is then given by $z := x + y$
\end{flushleft}
%
\begin{flushleft}
b) Since $x$ and $y$ are i.i.d and $x \sim \mathcal{N}(\mu_k^x, \Sigma_k^x)$ and $y \sim \mathcal{N}(\mu_k^y, \Sigma_k^y)$, the sum of two Gaussian distributions is again Gaussian. Therefore $p(z \,|\, \theta^x, \theta^y)$ is a mixture of Gaussians.
\end{flushleft}
%
\begin{flushleft}
c) 
\[ p(z \,|\, \theta^x, \theta^y) = \sum_{k=1}^{K_x} \sum_{l=1}^{K_y} \pi_k^x \pi_l^y \; \mathcal{N}(z \,|\, \mu_k^x + \mu_l^y, \Sigma_k^x + \Sigma_l^y) \]
\end{flushleft}
%
\Problem{5}
%
%
%

%%%%%%%%%%%%%%%%%%%%%%%%%%%%%%%%%%%%%%%%%%%%%%%%%%%%%%%%%%%%%%%%%%%%%%%%%%%%%%%%
% !!! HOMEWORK ENDS HERE !!!
%%%%%%%%%%%%%%%%%%%%%%%%%%%%%%%%%%%%%%%%%%%%%%%%%%%%%%%%%%%%%%%%%%%%%%%%%%%%%%%%

%%%%%%%%%%%%%%%%%%%%%%%%%%%%%%%%%%%%%%%%%%%%%%%%%%%%%%%%%%%%%%%%%%%%%%%%%%%%%%%%
\newpage

\vspace*{-15.8mm}
\fontsize{19pt}{21pt}\selectfont

\vspace{25.3mm}
Appendix

\normalsize\selectfont
\vspace{13.2mm}
We confirm that the submitted solution is original work and was written by us without further assistance. Appropriate credit has been given where reference has been made to the work of others.

\vspace{18.1mm}
\rule[-3.7mm]{\linewidth}{0.5pt}
Munich, \today, Signature \PersonOne

\vspace{18.1mm}
\rule[-3.7mm]{\linewidth}{0.5pt}
Munich, \today, Signature \PersonTwo

\vspace{18.1mm}
\rule[-3.7mm]{\linewidth}{0.5pt}
Munich, \today, Signature \PersonThree
 % !!! DON'T TOUCH !!!
%%%%%%%%%%%%%%%%%%%%%%%%%%%%%%%%%%%%%%%%%%%%%%%%%%%%%%%%%%%%%%%%%%%%%%%%%%%%%%%%

\end{document}
